%\documentclass[handout]{beamer}
\documentclass[ignorenonframetext]{beamer}
\usepackage{textpos}
\usepackage{graphicx}
\usepackage{pgf}
\usepackage{caption}
\usepackage{listings}
\usepackage{multimedia}
\usepackage{gensymb}
\usepackage{amsmath, mathrsfs}

% \captionsetup[figure]{labelformat=empty}% redefines the caption setup of the figures environment in the beamer class.

\lstdefinestyle{yaml}{
     basicstyle=\ttfamily\color{blue}\tiny,
     rulecolor=\color{black},
     string=[s]{'}{'},
     stringstyle=\color{blue},
     comment=[l]{:},
     commentstyle=\color{black},
     morecomment=[l]{-}
 }

\usetheme{Boadilla}
\usefonttheme{serif}
% \mode<presentation>
% {
%  \usefonttheme{serif}
% % \useoutertheme{sidebar}
% %   \logo{\includegraphics[height=1cm]{elec_logo.pdf}}
% }

\title[TART Imaging]{Imaging. A Visual Explanation (with TART)}

\author[Molteno]{Tim Molteno}

\institute[Otago]
{
  Electronics Research Foundation \\
  \& \\
  Department of Physics,
  University of Otago \\
  \vspace{1cm}
  \large{Dunedin, New Zealand.}\\
  \vspace{2cm}
  \includegraphics[width=0.6\linewidth]{../tart_overview/fig/elec_header_font.pdf}
}

\logo{\pgfputat{\pgfxy(-0.72,7.7)}{\pgfbox[center,base]{\includegraphics[width=1cm]{../tart_overview/fig/elec_logo.pdf}}}}

\date[NASSP 04/2025] % (optional, should be abbreviation of conference name)
{}

% \addtobeamertemplate{headline}{}{%
% \begin{textblock*}{100mm}(0.87\textwidth,2mm)
% \includegraphics[height=1.5cm]{elec_logo.pdf}
% \end{textblock*}}

\begin{document}

% Abstract: In this talk I'll give an overview of Africa's newest radio telescope, the TART, recently installed in Mauritius in a joint effort between SARAO, the University of Otago (NZ). I'll talk about TARTs origins, it's design and open-source philosophy as well as the upcoming project to install TART telescopes in the SKA African partner nations. I'll also include some details about improvements that will appear in TART-3, the next version of TART!

\begin{frame}
  \titlepage
\end{frame}
 
% \begin{frame}
% \vspace{1cm}
%   \includegraphics[width=\linewidth]{../tart_overview/fig/spiral.jpeg}\\
%   \centering{A TART radio telescope under construction in Mauritius}
% \end{frame}
%

\begin{frame}
  \tableofcontents
  % You might wish to add the option [pausesections]
\end{frame}


\section{Visibilities: What are they}

\frame{\tableofcontents[currentsection]}

\begin{frame}{Synthesis Imaging}

\begin{block}{What are visibilities?}
 Just the averages of the product of pairs of signals $E_i(t) \times E_j(t)$.
\end{block}
\pause
\begin{columns}
 \begin{column}{0.55\linewidth}
\begin{block}{Really}
  Yes, just multiplication and averaging
\[ V_{ij} = \frac{1}{T} \int_0^T E_i(t) E_j^{\star}(t) dt \]
\end{block}
\begin{enumerate}
 \item 276 Pairs of antennas
 \item T $\sim$ 1 second
 \item Real-time correlation in FPGA
 \item 4.5 Giga MAC per second
\end{enumerate}
 \end{column}
 \begin{column}{0.45\linewidth}
   \includegraphics[width=\linewidth]{../tart_overview/fig/papilio_pro.jpg}
 \end{column}
\end{columns}
\end{frame}


\begin{frame}
Point sources make plane waves. Received as voltages...
 \begin{eqnarray*}
E_i & = & A e^{i (k r - \omega t)}, \\
E_j & = & A e^{i (k (r + dr) - \omega (t+dt)} \\
    & = & A e^{ikr} e^{ik d r} e^{-i \omega t}e^{-i \omega d t} = E_i e^{i (k dr - \omega dt)}
 \end{eqnarray*}
 \begin{eqnarray*}
  V_{ij} & = & \frac{1}{T} \int_0^T E_i(t) E_j^{\star}(t) dt \\
   & = & \frac{1}{T} \int_0^T A e^{i (k r - \omega t)} E_i^{\star} e^{-i (k dr - \omega dt)} dt \\
   & = & \frac{1}{T} \int_0^T A e^{i (k r - \omega t)} A e^{-i (k r - \omega t)}  e^{-i (k dr - \omega dt)} \\
   & = & A^2 \frac{1}{T} \int_0^T e^{-i (k dr - \omega dt)} = A^2 e^{-i (k dr - \omega dt)}
 \end{eqnarray*}
 \begin{block}{Yep its that simple}
  Visibilities depend on the intensity of the source ($A^2$), and the phase difference (depends on angle and baseline length)
 \end{block}
\end{frame}


\begin{frame}
\begin{columns}
 \begin{column}{0.3\linewidth}
  Move a source around: plot the visibility.
  \begin{block}{Quiz}
   Where are the baselines?
  \end{block}
\end{column}
 \begin{column}{0.7\linewidth}
  \begin{center}
  \includegraphics[width=0.45\linewidth]{images/harmonic_uvw_0.pdf}
  \includegraphics[width=0.45\linewidth]{images/harmonic_uvw_1.pdf}\\
  \includegraphics[width=0.45\linewidth]{images/harmonic_uvw_2.pdf}
  \includegraphics[width=0.45\linewidth]{images/harmonic_uvw_3.pdf}
  \end{center}
 \end{column}
\end{columns}
\end{frame}


\begin{frame}{Visibility, Angles and Spacing}
Well its geometry. A bit of diagram drawing gives.
 \begin{eqnarray*} V_{i,j} & = & A^2 e^{2 i \pi \left(\Delta x \sin{\left(\phi_{0} \right)} \sin{\left(\theta_{0} \right)} + \Delta y \sin{\left(\theta_{0} \right)} \cos{\left(\phi_{0} \right)} + \Delta z \cos{\left(\theta_{0} \right)}\right)}  \\
 & = & A^2 F(\Delta x, \Delta y)(\theta_0, \phi_0)
 \end{eqnarray*}
 where:
 \begin{itemize}
  \item $\Delta x_{ij}$ is the distance between antenna in $x$
  \item $\Delta y_{ij}$ is the distance between antennas in $y$
  \item $\Delta z_{ij}$ is the distance between antennas in $z$ (usually zero)
  \item We know these and they don't change
 \end{itemize}
 \begin{block}{Astronomers}
  We don't like $\Delta x$,$\Delta y$ so call them $u, v$.
 \end{block}
\end{frame}


\begin{frame}
\begin{columns}
 \begin{column}{0.3\linewidth}
  \begin{block}{Fringes}
   One for each baseline?
  \end{block}
  \[ F_{ij}(\theta, \phi) \]
\end{column}
 \begin{column}{0.7\linewidth}
  \begin{center}
  \includegraphics[width=0.45\linewidth]{images/harmonic_uvw_0.pdf}
  \includegraphics[width=0.45\linewidth]{images/harmonic_uvw_1.pdf}\\
  \includegraphics[width=0.45\linewidth]{images/harmonic_uvw_2.pdf}
  \includegraphics[width=0.45\linewidth]{images/harmonic_uvw_3.pdf}
  \end{center}
 \end{column}
\end{columns}
\end{frame}

\begin{frame}{All TART Fringes}
  \includegraphics[width=\linewidth]{images/all_fringes.png}
\end{frame}

\section{Imaging is now easy}

\frame{\tableofcontents[currentsection]}

\begin{frame}{Ridiculously Simple}
 \[ S(\theta, \phi) = \sum_{i,j} V_{ij} F_{ij}(\theta, \phi)\]
 \begin{block}{Just add them up}
  The sky is the sum of the fringes multiplied by their visibility.
 \end{block}
\end{frame}

\begin{frame}{Adding fringes}
\begin{minipage}{\linewidth}
  \centering
  \raisebox{-0.4\height}{\includegraphics[height=2cm]{images/harmonics/fringe_0.png}} +
  \raisebox{-0.4\height}{\includegraphics[height=2cm]{images/harmonics/fringe_1.png}} =
  \raisebox{-0.4\height}{\includegraphics[height=2cm]{images/harmonics/sum_1.png}}
\end{minipage}
\begin{minipage}{\linewidth}
  \centering
  \raisebox{-0.4\height}{\includegraphics[height=2cm]{images/harmonics/fringe_0.png}} +
  \raisebox{-0.4\height}{\includegraphics[height=2cm]{images/harmonics/fringe_1.png}} +
  \raisebox{-0.4\height}{\includegraphics[height=2cm]{images/harmonics/fringe_2.png}} =
  \raisebox{-0.4\height}{\includegraphics[height=2cm]{images/harmonics/sum_2.png}}
\end{minipage}
\begin{minipage}{\linewidth}
  \centering
  \raisebox{-0.4\height}{\includegraphics[height=1cm]{images/harmonics/fringe_0.png}} +
  \raisebox{-0.4\height}{\includegraphics[height=1cm]{images/harmonics/fringe_1.png}} +
  \raisebox{-0.4\height}{\includegraphics[height=1cm]{images/harmonics/fringe_2.png}} +
  \raisebox{-0.4\height}{\includegraphics[height=1cm]{images/harmonics/fringe_3.png}} +
  \raisebox{-0.4\height}{\includegraphics[height=1cm]{images/harmonics/fringe_4.png}} +
  \raisebox{-0.4\height}{\includegraphics[height=1cm]{images/harmonics/fringe_5.png}} =
  \raisebox{-0.4\height}{\includegraphics[height=1cm]{images/harmonics/sum_5.png}}
\end{minipage}
\begin{minipage}{\linewidth}
  \centering
  \raisebox{-0.4\height}{\includegraphics[height=1cm]{images/harmonics/fringe_0.png}} +
  \raisebox{-0.4\height}{\includegraphics[height=1cm]{images/harmonics/fringe_1.png}} +
  \raisebox{-0.4\height}{\includegraphics[height=1cm]{images/harmonics/fringe_2.png}} +
  \raisebox{-0.4\height}{\includegraphics[height=1cm]{images/harmonics/fringe_3.png}} +
  \raisebox{-0.4\height}{\includegraphics[height=1cm]{images/harmonics/fringe_4.png}} \ldots
  \raisebox{-0.4\height}{\includegraphics[height=1cm]{images/harmonics/fringe_51.png}} =
  \raisebox{-0.4\height}{\includegraphics[height=1cm]{images/harmonics/sum_51.png}}
\end{minipage}
\begin{minipage}{\linewidth}
\centering
\includegraphics[height=2cm]{images/harmonics/sum_275.png}
\end{minipage}
\end{frame}

\begin{frame}{Imaging}
 Imaging: work out where the sources are from 276 measured visibilities. So solve 276 equations to work out the sky.
 \begin{itemize}
 \item If the sky has 5 unknown point sources, we are totally OK.
  \item If the sky is an image, we have 1 million unknown pixels.
  \item If we only have 276 equations and 1 million unknowns, we have to make assumptions.
 \end{itemize}
\begin{block}{Can we use least squares?}
 Yes, but we still have to use assumptions.
\end{block}
\begin{block}{What kind of assumptions?}
 We could pretend that there are millions of unmeasured visibilities. And that these are all zero.
\end{block}
\end{frame}

\begin{frame}{Lets try it}
\includegraphics[width=\linewidth]{images/demo-time.jpg}
\end{frame}


% \begin{frame}{Calibration}
%  This is dark magic. Ben and Oleg are the experts. This is where all the work goes in on real telescopes.
% \end{frame}
% 
% \begin{frame}{Synthesis Continued...}
% Fourier transform relationship between radio sky brightness $I(l,m)$ and
% visibility $V(u,v)$,
% \[
%   V(u,v) = \int I(l,m) e^{2\pi j(lu  + mv)} dl dm
% \]
% Obtain $I(l,m)$ through inverse Fourier transform of $V(u,v)$, 
% \[
%  I(l,m) = \mathscr{F}^{-1}\{V(u,v)\}
% \]
% where $V(u,v)$ is the visibility function sampled in the $uv$-plane at the locations of each antenna $(u_{ij},v_{ij})$ pair. 
% \end{frame}
% 

% \section{Expected \& Unexpected Results}


\section{Stimela pipeline for TART}


\begin{frame}{Stimela}
 Modern framework for automating radio astronomy imaging. Describes each steps required to process the data. In our case these are:
 \begin{enumerate}
  \item Download data from a TART (mu-udm)
  \item Create a measurement set
  \item Image the measurement set
  \item Write the image to a 2D image format
 \end{enumerate}
We describe this using a stimela recipe. Just a text file.
\begin{block}{Docs}
\url{https://stimela.readthedocs.io/en/latest/}
\end{block}
\end{frame}


\begin{frame}[fragile]{Stimela: Recipes for processing data}

\begin{lstlisting}[style=yaml]
_include:
    - (tartcargo):
        - tart2ms.yml
        - tart-download-data.yml
        - spotless.yml
        - disko-draw.yml

tart-image:
    inputs:
        tart:
            dtype: str
            required: true
            info: "TART name (e.g. mu-udm)"

    outputs:
        ms:
            dtype: MS
            default:  '{current.tart}.ms'
        hdf:
            dtype: File
            default: '{current.tart}-vis.hdf'
        svg:
            dtype: File
            default: '{current.tart}.svg'

    assign:
        api: 'https://api.elec.ac.nz/tart/{recipe.tart}'

    steps:
        download-hdf:
            cab: tart-download-data
            skip: =EXISTS(recipe.hdf)
            params:
                api: =recipe.api
                vis: true
                file: =recipe.hdf

        create-ms:
            cab: tart2ms
            skip: =EXISTS(recipe.ms)
            params:
                hdf: =steps.download-hdf.file
                ms: =recipe.ms
                clobber: true
                single-field: true
                rephase: obs-midpoint

        spotless-image:
            cab: spotless
            params:
                ms: =recipe.ms
                HDF: ='{recipe.tart}-spotless.sphere'
                nvis: 10000
                multimodel: true
                healpix: true
                fov: 170deg
                res: 0.5deg

        spotless-draw:
            cab: disko-draw
            params:
                hdf: ='{recipe.tart}-spotless.sphere'
                show-sources: true
                scale-mad: true
                SVG: ='spotless-{recipe.svg}'
\end{lstlisting}
\end{frame}
\subsection{TART Cargo}

\begin{frame}[fragile]{TART Cargo}
%    \includegraphics[width=\linewidth]{fig/albatross_frame.png}
  The TART software for imaging run in an environment called \href{https://github.com/caracal-pipeline/stimela}{stimela}. This is the software framework for the worlds largest radio telescopes.

  \begin{block}{Install}
  \begin{verbatim}
   pip install tart-cargo
  \end{verbatim}
  \end{block}

  Notes:
  \begin{itemize}
   \item Works on Linux and Windows (WSL2).
   \item On a mac? Run a linux virtual machine.
  \end{itemize}

\end{frame}


\begin{frame}[fragile]{Sample Imaging Recipe}
There is a sample recipe that should also be run:
\begin{itemize}
 \item Download it from \url{https://github.com/tart-telescope/tart_cargo}
 \item In \verb!test/example-recipe.yml!
\end{itemize}
\begin{block}{Run It}
 \begin{verbatim}
  stimela run example_recipe.yml tart=tart-kenya
 \end{verbatim}
\end{block}
 The first time this is run, it will take a while and download a ~1GB docker image.
\end{frame}

\begin{frame}{Output}
\begin{columns}
 \begin{column}{0.7\linewidth}
\includegraphics[width=\linewidth]{images/obs_00000.hdf.png}
 \end{column}
 \begin{column}{0.3\linewidth}
 \begin{itemize}
  \item Known sources circled.
  \item Round!
 \end{itemize}
 \end{column}
\end{columns}
\end{frame}


\subsection{What is going on behind the scenes}

\frame{\tableofcontents[currentsection]}

\begin{frame}
 \includegraphics[width=\linewidth]{../tart_overview/fig/browser_view.png}
\end{frame}

\begin{frame}[containsverbatim]
\frametitle{TART API: Access data from anywhere}
TART data and control are done throught a web-based API.

\centering{\url{https://api.elec.ac.nz/tart/tart-kenya/}}
%
% \begin{lstlisting}[language=Python, frame=single, basicstyle=\footnotesize]
%  import requests
%
%  api_endpoint = "https://api.elec.ac.nz/tart/bw-biust/api/"
%  r = requests.get(api_endpoint + "/v1/imaging/vis")
%  print(r.json())
% \end{lstlisting}

\end{frame}

\begin{frame}{Conclusions}
 \begin{columns}
  \begin{column}{0.6\linewidth}
    \begin{itemize}
    \item Imaging isn't too hard.
    \item Projects ideas: Trying different TART imaging ideas
    \item Project idea: Use Machine learning!
    \item More projects: Tools for calibration of arrays.
    \item More projects: Bigger, more directional antennas.
    \end{itemize}
  \end{column}
  \begin{column}{0.4\linewidth}
    \includegraphics[width=\linewidth]{../tart_overview/fig/tart3.jpg}
  \end{column}
\end{columns}
    \begin{block}{}
     \begin{center} Questions? \end{center}
    \end{block}

\end{frame}



\end{document}
